\documentclass[12pt]{amsart}

%==================== Packages ====================
\usepackage[margin=1in]{geometry}
\usepackage{amsmath,amssymb,amsthm}
\usepackage{mathtools}              % paired delimiters, \coloneqq, etc.
\usepackage{enumitem}               % control list labels
\usepackage{bm}                     % bold math: \bm{v}
\usepackage{xcolor}
\usepackage[colorlinks=true,linkcolor=blue,urlcolor=blue]{hyperref}  % load last

\allowdisplaybreaks                 % let long align blocks break across pages
\setlist[enumerate,1]{label=(\alph*)}  % problem parts as (a), (b), (c)

%==================== Number systems ====================
\newcommand{\N}{\mathbb{N}}
\newcommand{\Z}{\mathbb{Z}}
\newcommand{\Q}{\mathbb{Q}}
\newcommand{\R}{\mathbb{R}}
\newcommand{\C}{\mathbb{C}}
\newcommand{\F}{\mathbb{F}}

%==================== Linear algebra ====================
\DeclareMathOperator{\Ran}{Ran}         % range (Treil); \ker already exists
\DeclareMathOperator{\rank}{rank}
\DeclareMathOperator{\nullity}{nullity}
\DeclareMathOperator{\spn}{span}        % NOT \span: that's a TeX primitive
\DeclareMathOperator{\tr}{tr}
\DeclareMathOperator{\col}{Col}
\DeclareMathOperator{\row}{Row}
\DeclareMathOperator{\proj}{proj}
\DeclareMathOperator{\id}{id}
\newcommand{\Mat}[2]{M_{#1\times #2}(\F)}   % \Mat{m}{n} -> M_{m x n}(F)
\newcommand{\Pn}[1][n]{\mathcal{P}_{#1}}    % \Pn -> P_n, \Pn[3] -> P_3
\newcommand{\tp}{^{\mathsf{T}}}             % A\tp -> A^T
\newcommand{\inv}{^{-1}}                    % A\inv -> A^{-1}
\newcommand{\mat}[1]{\begin{bmatrix}#1\end{bmatrix}}
\newcommand{\pmat}[1]{\begin{pmatrix}#1\end{pmatrix}}

%==================== Analysis / topology ====================
\DeclareMathOperator{\diam}{diam}
\DeclareMathOperator{\dist}{dist}
\DeclareMathOperator{\interior}{int}    % NOT \int: that's the integral sign
\DeclareMathOperator{\sgn}{sgn}
\DeclareMathOperator{\supp}{supp}
\newcommand{\closure}[1]{\overline{#1}}
\newcommand{\pre}{^{\mathrm{pre}}}          % f\pre(V), Pugh's preimage
\newcommand{\nbhd}[2]{N_{#1}(#2)}           % \nbhd{r}{p} -> N_r(p), Rudin
\newcommand{\limn}{\lim_{n\to\infty}}
\newcommand{\eps}{\varepsilon}
\newcommand{\dd}{\,\mathrm{d}}              % \int f(x)\dd x
\newcommand{\restr}[2]{\left.#1\right|_{#2}}  % \restr{f}{A} -> f|_A
\newcommand{\seq}[2][n]{\left(#2\right)_{#1\in\N}}  % \seq{p_n} -> (p_n)_{n in N}

%==================== Delimiters (mathtools) ====================
% \abs{x} normal size; \abs*{x} auto-sizes; \abs[\big]{x} manual size
\DeclarePairedDelimiter{\abs}{\lvert}{\rvert}
\DeclarePairedDelimiter{\norm}{\lVert}{\rVert}
\DeclarePairedDelimiter{\floor}{\lfloor}{\rfloor}
\DeclarePairedDelimiter{\ceil}{\lceil}{\rceil}
\DeclarePairedDelimiter{\inner}{\langle}{\rangle}
\DeclarePairedDelimiterX{\set}[1]{\lbrace}{\rbrace}{#1}
\DeclarePairedDelimiterX{\setof}[2]{\lbrace}{\rbrace}{#1 \,:\, #2}  % \setof{x\in\R}{x>0}

%==================== Logic shortcuts ====================
\newcommand{\contra}{\Rightarrow\!\Leftarrow}   % contradiction marker
\newcommand{\st}{\text{ such that }}

%==================== Theorem environments ====================
\theoremstyle{plain}
\newtheorem*{theorem*}{Theorem}
\newtheorem*{lemma*}{Lemma}
\newtheorem*{claim*}{Claim}
\theoremstyle{definition}
\newtheorem*{definition*}{Definition}
\theoremstyle{remark}
\newtheorem*{remark*}{Remark}

\renewcommand{\qedsymbol}{$\blacksquare$}       % black box

% Problem header using the textbook's exercise number: \begin{problem}{2.5}
\newenvironment{problem}[1]{\par\medskip\noindent\textbf{#1.}\ }{\par\smallskip}
% For computational answers that aren't proofs
\newenvironment{solution}{\begin{proof}[Solution]}{\end{proof}}

\title{Homework 7}
\author{Antony Perez}

\begin{document}

\maketitle

%--------------------------------------------------
\begin{problem}{3.1}
If $A$ is an $n\times n$ matrix, how are the determinants $\det A$ and $\det(5A)$ related?

\noindent\textbf{Remark:} $\det(5A) = 5\det A$ only in the trivial case of $1\times 1$ matrices.
\end{problem}

\begin{solution}
Let $\det (A) = \alpha$. These two are related because $\det (5A) = 5^n \times \alpha$. This is because we can pull out the scalar 5 out of each column vector, and, since there are $n$ column vectors, we do this $n$ times, giving us $5^n \det (A) = 5^n \times \alpha$.
\end{solution}

%--------------------------------------------------
\begin{problem}{3.2}
How are the determinants $\det A$ and $\det B$ related if
\begin{enumerate}
  \item
  \[
    A = \pmat{a_1 & a_2 & a_3 \\ b_1 & b_2 & b_3 \\ c_1 & c_2 & c_3}, \qquad
    B = \pmat{2a_1 & 3a_2 & 5a_3 \\ 2b_1 & 3b_2 & 5b_3 \\ 2c_1 & 3c_2 & 5c_3};
  \]
  \item
  \[
    A = \pmat{a_1 & a_2 & a_3 \\ b_1 & b_2 & b_3 \\ c_1 & c_2 & c_3}, \qquad
    B = \pmat{3a_1 & 4a_2 + 5a_1 & 5a_3 \\ 3b_1 & 4b_2 + 5b_1 & 5b_3 \\ 3c_1 & 4c_2 + 5c_1 & 5c_3}.
  \]
\end{enumerate}
\end{problem}

\begin{solution}
    For $(a)$, they are related by linearity. For $(b)$, they are also related by linearity and row/column addition. Pulling the scalar out of each column gives $\det B = 2\cdot 3\cdot 5 \det A = 30\det A$ in (a). In (b), splitting column 2 gives $3\cdot 4\cdot 5\det A$ plus a determinant with two proportional columns, which is 0, so $\det B = 60\det A$.
\end{solution}

%--------------------------------------------------
\begin{problem}{3.3}
Using column or row operations compute the determinants
\[
  \begin{vmatrix} 0 & 1 & 2 \\ -1 & 0 & -3 \\ 2 & 3 & 0 \end{vmatrix}, \quad
  \begin{vmatrix} 1 & 2 & 3 \\ 4 & 5 & 6 \\ 7 & 8 & 9 \end{vmatrix}, \quad
  \begin{vmatrix} 1 & 0 & -2 & 3 \\ -3 & 1 & 1 & 2 \\ 0 & 4 & -1 & 1 \\ 2 & 3 & 0 & 1 \end{vmatrix}, \quad
  \begin{vmatrix} 1 & x \\ 1 & y \end{vmatrix}.
\]
\end{problem}

\begin{solution}
    Let $A,B,C,D$ be the matrices above respectively. For $A$, swap $R_1$ with $R_2$, which flips the sign. Then $R_3+ 2R_1$. We finally have $r_3 - 3R_2$ We now multiply across the diagonal and get $-((-1)(1)(-12))= -12$. $\det(A) = -12$. \\
    For $B$, $R_2-4R_1; R_3 - 7R_1; R_3 - 2R_2.$ This makes Row 3 all zeros, showing linear dependence, which shows $\det(B)=0.$\\
    For $C$, $R_2 + 3R_1; R_4 - 2R_1; R_3-4R_2;  R_4 - 3R_2; R_4 - R_3$. This gives us a triangular matrix, so $\det(C) = 1 \cdot 1\cdot19\cdot 5 = 95$. \\
    For $D$, $R_2 - R_1$. $\det (D)=y-x$.
\end{solution}
%--------------------------------------------------
\begin{problem}{3.4}
A square ($n\times n$) matrix is called \emph{skew-symmetric} (or \emph{antisymmetric}) if $A^T = -A$. Prove that if $A$ is skew-symmetric and $n$ is odd, then $\det A = 0$. Is this true for even $n$?
\end{problem}

\begin{proof}
 
Since $A^T=-A$, taking determinants gives
\[
  \det(A)=\det(A^T)=\det(-A)=(-1)^n\det(A).
\]
Since $n$ is odd, $\det(A)=-\det(A)$, so $\det(A)=0$.
This is not true for even $n$.
$A=\pmat{0&1\\-1&0}$
is skew-symmetric, but $\det(A)=1$.
\end{proof}

%--------------------------------------------------
\begin{problem}{3.5}
A square matrix is called \emph{nilpotent} if $A^k = \bm{0}$ for some positive integer $k$. Show that for a nilpotent matrix $A$, $\det A = 0$.
\end{problem}

\begin{proof}
Since $A$ is nilpotent, there exists $k \in \N$ such that $A^k = \bm{0}$.
By Theorem 3.5, we get
\[
  (\det A)^k = \det(A^k) = \det(\bm{0}) = 0.
\]
This implies $\det A = 0$.
\end{proof}

%--------------------------------------------------
\begin{problem}{3.6}
Prove that if the matrices $A$ and $B$ are similar, then $\det A = \det B$.
\end{problem}

\begin{proof}
  Let $A = PBP^{-1}$. $\det(A) = \det(PBP^{-1}) \to \det(PP^{-1})\det(B) \implies \det A = \det B$. 
\end{proof}

%--------------------------------------------------
\begin{problem}{3.7}
A real square matrix $Q$ is called \emph{orthogonal} if $Q^T Q = I$. Prove that if $Q$ is an orthogonal matrix then $\det Q = \pm 1$.
\end{problem}

\begin{proof}
  Let $Q^TQ = I$. $\det(Q^TQ)=\det(I) \implies \det(Q^T)\det(Q)=1$. Since $\det(Q^T) = \det(Q)$, they must have the same value. Since they're being multiplied together, the only solutions to this equation is if $\det(Q^T)=\det(Q)=\pm 1$. This concludes the proof. 
\end{proof}

%--------------------------------------------------
\begin{problem}{3.8}
Show that
\[
  \begin{vmatrix} 1 & x & x^2 \\ 1 & y & y^2 \\ 1 & z & z^2 \end{vmatrix}
  = (z - x)(z - y)(y - x).
\]
This is a particular case of the so-called Vandermonde determinant.
\end{problem}

\begin{proof}
Subtracting the first row from the second and third rows gives
\[
\det(A)=
\begin{vmatrix}
1&x&x^2\\
0&y-x&y^2-x^2\\
0&z-x&z^2-x^2
\end{vmatrix}.
\]
Factoring out $y-x$ and $z-x$, then subtracting the second row from the third, gives
\[
\det(A)=(y-x)(z-x)
\begin{vmatrix}
1&x&x^2\\
0&1&y+x\\
0&0&z-y
\end{vmatrix}.
\]
Multiplying across the diagonal, we get
\[
\det(A)=(z-x)(z-y)(y-x).
\]
\end{proof}

%--------------------------------------------------
\begin{problem}{3.9}
Let points $A$, $B$ and $C$ in the plane $\R^2$ have coordinates $(x_1, y_1)$, $(x_2, y_2)$ and $(x_3, y_3)$ respectively. Show that the area of triangle $ABC$ is the absolute value of
\[
  \frac{1}{2}
  \begin{vmatrix} 1 & x_1 & y_1 \\ 1 & x_2 & y_2 \\ 1 & x_3 & y_3 \end{vmatrix}.
\]

\noindent\textbf{Hint:} use row operations and geometric interpretation of $2\times 2$ determinants (area).
\end{problem}

\begin{proof}
Subtracting the first row from the second and third rows gives
\[
\det(A)=
\begin{vmatrix}
1&x_1&y_1\\
0&x_2-x_1&y_2-y_1\\
0&x_3-x_1&y_3-y_1
\end{vmatrix}
=
\begin{vmatrix}
x_2-x_1&y_2-y_1\\
x_3-x_1&y_3-y_1
\end{vmatrix}.
\]
The absolute value of this determinant gives the area of the parallelogram determined by the vectors from $A$ to $B$ and $C$. Since the triangle has half this area,
\[
\operatorname{Area}(ABC)=\frac12\left|
\begin{vmatrix}
1&x_1&y_1\\
1&x_2&y_2\\
1&x_3&y_3
\end{vmatrix}
\right|.
\]
\end{proof}

%--------------------------------------------------
\begin{problem}{3.10}
Let $A$ be a square matrix. Show that block triangular matrices
\[
  \pmat{I & * \\ \bm{0} & A}, \qquad
  \pmat{A & * \\ \bm{0} & I}, \qquad
  \pmat{I & \bm{0} \\ * & A}, \qquad
  \pmat{A & \bm{0} \\ * & I}
\]
all have determinant equal to $\det A$. Here $*$ can be anything.
\end{problem}

\begin{solution}
    All four matrices are block triangular, so their determinants equal the products of the determinants of their diagonal blocks. Since $\det(I)=1$, each determinant is $\det(A)$.
\end{solution}

\medskip
The following problems illustrate the power of block matrix notation.

%--------------------------------------------------
\begin{problem}{3.11}
Use the previous problem to show that if $A$ and $C$ are square matrices, then
\[
  \det \pmat{A & B \\ \bm{0} & C} = \det A \det C.
\]

\noindent\textbf{Hint:}
$\pmat{A & B \\ \bm{0} & C} = \pmat{I & B \\ \bm{0} & C} \pmat{A & \bm{0} \\ \bm{0} & I}$.
\end{problem}

\begin{proof}
Using the factorization in the hint and taking determinants, we get
\[
\det\pmat{A&B\\\bm{0}&C}
=
\det\pmat{I&B\\\bm{0}&C}
\det\pmat{A&\bm{0}\\\bm{0}&I}.
\]
By Problem 3.10, these determinants are $\det(C)$ and $\det(A)$, respectively. Therefore,
\[
\det\pmat{A&B\\\bm{0}&C}=\det(A)\det(C).
\]
\end{proof}


%--------------------------------------------------
\begin{problem}{3.12}
Let $A$ be $m\times n$ and $B$ be $n\times m$ matrices. Prove that
\[
  \det \pmat{0 & A \\ -B & I} = \det(AB).
\]

\noindent\textbf{Hint:} While it is possible to transform the matrix by row operations to a form where the determinant is easy to compute, the easiest way is to right multiply the matrix by $\pmat{I & 0 \\ B & I}$.
\end{problem}

\begin{proof}
Right multiplying by the matrix in the hint gives
\[
\pmat{0&A\\-B&I}
\pmat{I&0\\B&I}
=
\pmat{AB&A\\0&I}.
\]
Taking determinants and applying Problems 3.10 and Theorem 3.5, we get
\[
\det\pmat{0&A\\-B&I}
=
\det(AB),
\]
since $\det\pmat{I&0\\B&I}=1$ and
$\det\pmat{AB&A\\0&I}=\det(AB)$.
\end{proof}

\end{document}